% ECG Reading Guide for Nurses - 護理心電圖判讀教學海報
% 作者: 謝慕揚, MD, PhD, FESC
% 資料來源: The Visual Nurse (@thevisualnurse)
\documentclass[12pt,a4paper]{article}
% Drake_preamble.tex - LaTeX Preamble Template
% 作者: 謝慕揚, MD, PhD, FESC
% 
% 使用說明:
% 1. 表格請使用 longtable，不要有垂直分隔線 |
% 2. 表格內換行使用 \newline，不要使用 \\
% 3. 藥物名稱、檢驗、檢查請維持英文
% 4. Reference 格式請使用 \href 加上驗證過的 DOI

% Including essential packages for Chinese typesetting
\usepackage{xeCJK}
\usepackage{fontspec}
% Configuring Chinese font with Noto Serif CJK TC, fallback to SimSun
\IfFontExistsTF{Noto Serif CJK TC}{
  \setCJKmainfont{Noto Serif CJK TC}
  \setCJKsansfont{Noto Serif CJK TC}
  \setCJKmonofont{Noto Serif CJK TC}
}{\setCJKmainfont{SimSun}
  \setCJKsansfont{SimSun}
  \setCJKmonofont{SimSun}
}

% Including other necessary packages
\usepackage{geometry}
\usepackage{titlesec}
\usepackage{enumitem}
\usepackage{xcolor}
\usepackage{colortbl}
\usepackage{tcolorbox}
\usepackage{booktabs}
\usepackage{longtable}
\usepackage{array}
\usepackage{graphicx}
\usepackage{tikz}
\usepackage{fancyhdr}
\usepackage{multicol}
\usepackage{datetime2}
\usepackage{hyperref}
\usepackage{mdframed}
\usepackage{amsmath}
\usepackage{amssymb}
\usepackage{multirow}

\hypersetup{
    colorlinks=true,
    linkcolor=emergency_blue,
    citecolor=emergency_blue,
    urlcolor=emergency_blue,
    pdfborder={0 0 0}
}

% Defining page geometry
\geometry{
  top=2.5cm,
  bottom=2.5cm,
  left=2cm,
  right=2cm,
  headheight=25pt
}

% Adjusting topmargin to compensate for increased headheight
\addtolength{\topmargin}{-10pt}

% Defining colors
\definecolor{emergency_red}{RGB}{186,24,27}
\definecolor{emergency_blue}{RGB}{0,114,188}
\definecolor{emergency_gray}{RGB}{120,120,120}
\definecolor{highlight_yellow}{RGB}{255,250,205}

% Setting title formats
\titleformat{\section}[block]{\Large\bfseries\color{emergency_red}}{\thesection}{10pt}{}
\titleformat{\subsection}[block]{\large\bfseries\color{emergency_blue}}{\thesubsection}{8pt}{}
\titleformat{\subsubsection}[block]{\normalsize\bfseries\color{emergency_gray}}{\thesubsubsection}{6pt}{}

% Defining custom environment for important notes
\newmdenv[
    linecolor=emergency_red,
    backgroundcolor=highlight_yellow,
    linewidth=2pt,
    topline=true,
    bottomline=true,
    leftline=true,
    rightline=true,
    innertopmargin=10pt,
    innerbottommargin=10pt,
    innerleftmargin=10pt,
    innerrightmargin=10pt
]{importantbox}

% Defining note command with tcolorbox
\newcommand{\note}[1]{
    \par
    \noindent
    \begin{tcolorbox}[
        colback=emergency_blue!5!white,
        colframe=emergency_blue,
        title=\textbf{重點提醒},
        fonttitle=\bfseries,
        boxrule=0.5pt,
        arc=3pt,
        left=6pt,
        right=6pt,
        top=6pt,
        bottom=6pt
    ]
    #1
    \end{tcolorbox}
}

% Key findings box
\newcommand{\keyfinding}[1]{
    \par
    \noindent
    \begin{tcolorbox}[
        colback=yellow!10!white,
        colframe=orange!75!black,
        title=\textbf{核心發現摘要},
        fonttitle=\bfseries,
        boxrule=0.5pt,
        arc=3pt
    ]
    #1
    \end{tcolorbox}
}

% Defining custom date format for ISO 8601
\DTMnewdatestyle{iso}{
  \renewcommand{\DTMdisplaydate}[4]{\number##1-\number##2-\number##3}
  \renewcommand{\DTMDisplaydate}[4]{\number##1-\number##2-\number##3}
}
\DTMsetdatestyle{iso}
\newcommand{\mydate}{\DTMtoday}


% Custom header for this document
\pagestyle{fancy}
\fancyhf{}
\fancyhead[L]{\textcolor{emergency_red}{心電圖判讀教學}}
\fancyhead[R]{\textcolor{emergency_blue}{ECG Reading Guide}}
\fancyfoot[C]{\textcolor{emergency_gray}{\thepage}}
\renewcommand{\headrulewidth}{1pt}
\renewcommand{\headrule}{\hbox to\headwidth{\color{emergency_red}\leaders\hrule height \headrulewidth\hfill}}

\begin{document}

% Title Page
\begin{center}
{\LARGE\bfseries\textcolor{emergency_red}{護理心電圖判讀教學指南}}\\[0.3cm]
{\Large\textbf{How to Read an ECG - Single Lead Telemetry}}\\[0.5cm]
{\large 病房護理人員快速參考手冊}\\[0.3cm]
{\normalsize 資料來源：The Visual Nurse (@thevisualnurse)}\\[0.3cm]
{\normalsize 整理：謝慕揚 MD, PhD, FESC \quad 日期：\mydate}
\end{center}

\vspace{0.5cm}

\begin{importantbox}
\textbf{心電圖判讀八步驟口訣：R-P-Q-STaT}
\begin{itemize}[leftmargin=*]
    \item \textbf{R} $\times 2$：Rate (心率) \& Rhythm (節律)
    \item \textbf{P} $\times 2$：P waves (P波) \& PR interval (PR間期)
    \item \textbf{Q} $\times 2$：QRS complex (QRS波群) \& QT interval (QT間期)
    \item \textbf{STaT}：ST segment (ST段) \& T waves (T波)
\end{itemize}
\end{importantbox}

%============================================================
\section{第一部分：心電圖基礎概念 (ECG Basics)}
%============================================================

\begin{center}
\includegraphics[width=\textwidth]{2026-01-11/ECG_01_Overview.png}
\end{center}

\subsection{心臟週期 (Cardiac Cycle)}

\begin{longtable}{p{4cm}p{10cm}}
\toprule
\textbf{階段} & \textbf{說明} \\
\midrule
\endfirsthead
\toprule
\textbf{階段} & \textbf{說明} \\
\midrule
\endhead
\bottomrule
\endfoot
心房去極化\newline (Atrial depolarization) & 產生 P 波，代表心房收縮前的電位變化 \\
\midrule
心室去極化\newline (Ventricular depolarization) & 產生 QRS 波群，代表心室收縮前的電位變化 \\
\midrule
心室再極化\newline (Ventricular repolarization) & 產生 T 波，代表心室恢復極化狀態 \\
\bottomrule
\end{longtable}

\subsection{心臟傳導系統 (Cardiac Conduction System)}

\begin{longtable}{p{1cm}p{4cm}p{9cm}}
\toprule
\textbf{順序} & \textbf{結構} & \textbf{功能} \\
\midrule
\endfirsthead
\toprule
\textbf{順序} & \textbf{結構} & \textbf{功能} \\
\midrule
\endhead
\bottomrule
\endfoot
1 & SA node\newline (竇房結) & 心臟的天然節律器，固有頻率 60-100 bpm \\
\midrule
2 & AV node\newline (房室結) & 延遲傳導，讓心房完成收縮後心室再收縮\newline 固有頻率 40-60 bpm \\
\midrule
3 & Bundle of His\newline (希氏束) & 傳導電位至心室 \\
\midrule
4 & Bundle branches\newline (束支) & 左右束支分別傳導至左右心室 \\
\midrule
5 & Purkinje fibers\newline (乏乾乏氏纖維) & 快速傳導至心室肌細胞\newline 固有頻率 20-40 bpm \\
\bottomrule
\end{longtable}

\subsection{正常數值記憶口訣}

\note{
\begin{center}
\textbf{0.12 to 0.20 (PR interval)} + \textbf{$<$0.12 (QRS)} = \textbf{$<$0.44 (QT interval)}
\end{center}
所有數值單位皆為「秒」(seconds)
}

\subsection{正常竇性心律特徵 (Normal Sinus Rhythm)}
\begin{itemize}[leftmargin=*]
    \item 每個 QRS 前有一個直立的 P 波 (One upright P wave per QRS)
    \item 所有 P 波外觀相似 (All P waves similar in appearance)
    \item PR interval 介於 0.12-0.20 秒
    \item 心率介於 60-100 bpm
\end{itemize}

\newpage
%============================================================
\section{第二部分：心率與節律分析 (Rate \& Rhythm Analysis)}
%============================================================

\begin{center}
\includegraphics[width=\textwidth]{2026-01-11/ECG_02_Rate_Rhythm_Pwaves.png}
\end{center}

\subsection{步驟 1：心率計算 (Heart Rate Calculation)}

\subsubsection{方法一：六秒法 (6-Second Method)}
\begin{itemize}[leftmargin=*]
    \item 數 6 秒內（30 大格）的 R 波數量
    \item R 波數量 $\times$ 10 = 每分鐘心率 (HR in BPM)
    \item \textbf{適用於}：不規則心律
\end{itemize}

\subsubsection{方法二：小格法 (Small Box Method)}
\begin{itemize}[leftmargin=*]
    \item 數兩個 R 波之間的小格數（每小格 = 0.04 秒）
    \item 1500 $\div$ 小格數 = 每分鐘心率
    \item \textbf{適用於}：規則心律
\end{itemize}

\subsubsection{方法三：大格法 (Big Box Method)}
\begin{itemize}[leftmargin=*]
    \item 找一個落在粗線上的 R 波
    \item 從該 R 波開始，數下一個 R 波經過幾條粗線
    \item 口訣：\textbf{300, 150, 100, 75, 60, 50}
    \item \textbf{適用於}：規則心律，快速估算
\end{itemize}

\subsection{步驟 2：節律規則性 (Rhythm Regularity)}

\keyfinding{
問自己：「P 波和 QRS 波群是否規律地排列？」\\
(Do the P waves and QRS complexes march out routinely?)
}

\begin{longtable}{p{6.5cm}p{7cm}}
\toprule
\textbf{規則心室反應}\newline (Regular ventricular response) & \textbf{不規則心室反應}\newline (Irregular ventricular response) \\
\midrule
\endfirsthead
R-R 間距相等 & R-R 間距不等 \\
可見於：Sinus rhythm, SVT, VT & 可見於：AF, MAT, Sinus arrhythmia \\
\bottomrule
\end{longtable}

\subsection{步驟 3：P 波分析 (P Wave Analysis)}

\begin{longtable}{p{4cm}p{4.5cm}p{5cm}}
\toprule
\textbf{類型} & \textbf{特徵} & \textbf{臨床意義} \\
\midrule
\endfirsthead
\toprule
\textbf{類型} & \textbf{特徵} & \textbf{臨床意義} \\
\midrule
\endhead
\bottomrule
\endfoot
Sinus rhythm\newline (竇性心律) & P 波直立且外觀相似\newline (Upright \& similar) & 正常傳導 \\
\midrule
Wandering atrial pacemaker\newline (游走性心房節律器) & P 波外觀交替變化\newline (Alternating appearance) & 心房內多個起搏點 \\
\midrule
Atrial fibrillation\newline (心房顫動) & 無明顯 P 波\newline (None present) & 心房不規則顫動，\newline 需注意中風風險 \\
\bottomrule
\end{longtable}

\newpage
%============================================================
\section{第三部分：間期與波段分析 (Intervals \& Segments)}
%============================================================

\begin{center}
\includegraphics[width=\textwidth]{2026-01-11/ECG_03_PR_QRS_QT_ST_Twaves.png}
\end{center}

\subsection{步驟 4：PR Interval (PR 間期)}

\begin{longtable}{p{4cm}p{10cm}}
\toprule
\textbf{數值} & \textbf{臨床意義} \\
\midrule
\endfirsthead
\bottomrule
\endfoot
正常：0.12-0.20 秒 & 正常房室傳導 \\
\midrule
$<$0.12 秒\newline + 倒置 P 波 & 可能為 Junctional rhythm (交界性心律)\newline 電位由房室結逆向傳導 \\
\midrule
$>$0.20 秒 & 可能為 AV block (房室傳導阻滯)\newline 或 AV delay (房室傳導延遲) \\
\bottomrule
\end{longtable}

\subsection{步驟 5：QRS Complex (QRS 波群)}

\begin{longtable}{p{5cm}p{8.5cm}}
\toprule
\textbf{QRS 寬度} & \textbf{臨床意義} \\
\midrule
\endfirsthead
\bottomrule
\endfoot
Narrow (窄) $<$0.12 秒 & 電位起源於心室以上 (Supraventricular origin)\newline 正常傳導路徑 \\
\midrule
Wide (寬) $>$0.12 秒 & 電位可能起源於心室 (Ventricular origin)\newline 或傳導異常 (Bundle branch block) \\
\bottomrule
\end{longtable}

\subsection{步驟 6：QT Interval (QT 間期)}

\note{
\textbf{正常值}：$<$0.44 秒，或小於 R-R interval 的一半\\
\textbf{測量方法}：從 QRS 起始到 T 波結束\\
\textbf{臨床重要性}：QT 延長會增加 Torsades de Pointes (TdP) 風險
}

\subsection{步驟 7：ST Segment (ST 段)}

評估 ST 段相對於等電位基線 (isoelectric baseline) 的偏移：

\begin{longtable}{p{4cm}p{5cm}p{4.5cm}}
\toprule
\textbf{變化} & \textbf{定義} & \textbf{可能原因} \\
\midrule
\endfirsthead
\bottomrule
\endfoot
ST elevation\newline (ST 上升) & ST 段高於基線 & \textcolor{emergency_red}{\textbf{STEMI}}、Pericarditis、\newline Early repolarization \\
\midrule
ST depression\newline (ST 下降) & ST 段低於基線 & Ischemia、NSTEMI、\newline Digoxin effect \\
\bottomrule
\end{longtable}

\subsection{步驟 8：T Wave (T 波)}

\begin{longtable}{p{3cm}p{5cm}p{5.5cm}}
\toprule
\textbf{型態} & \textbf{外觀} & \textbf{可能原因} \\
\midrule
\endfirsthead
\bottomrule
\endfoot
Normal (正常) & 直立、圓滑 & 正常再極化 \\
\midrule
Peaked (尖峰) & 高聳、尖銳 & \textcolor{emergency_red}{Hyperkalemia (高血鉀)} \\
\midrule
Inverted (倒置) & 向下 & Ischemia、MI、\newline Ventricular hypertrophy \\
\midrule
Biphasic (雙相) & 先上後下或先下後上 & Ischemia、Hypokalemia \\
\bottomrule
\end{longtable}

\newpage
%============================================================
\section{第四部分：心率參考值與練習單 (Reference \& Practice)}
%============================================================

\begin{center}
\includegraphics[width=\textwidth]{2026-01-11/ECG_04_Practice_Worksheet.png}
\end{center}

\subsection{各類心律固有頻率 (Intrinsic Rates)}

\begin{longtable}{p{4cm}p{3cm}p{3cm}p{3cm}}
\toprule
\textbf{起源} & \textbf{Bradycardia} & \textbf{Intrinsic Rate} & \textbf{Tachycardia} \\
\midrule
\endfirsthead
\bottomrule
\endfoot
\textbf{Sinus}\newline (竇房結) & $<$60 bpm & 60-100 bpm & $>$100 bpm \\
\midrule
\textbf{Junctional}\newline (交界性) & $<$40 bpm & 40-60 bpm & $>$100 bpm \\
\midrule
\textbf{Ventricular}\newline (心室) & -- & 20-40 bpm & 50-100 bpm (accelerated)\newline $>$100 bpm (VT) \\
\bottomrule
\end{longtable}

\note{
\textbf{SVT (Supraventricular tachycardia)}：依據發作機轉分類，心率通常 $>$150 bpm
}

\subsection{心律判讀練習清單}

每次判讀心電圖時，請依序回答以下問題：

\begin{enumerate}[leftmargin=*]
    \item \textbf{Atrial Rate (心房率)}：\_\_\_\_\_\_ bpm \quad \textbf{Ventricular Rate (心室率)}：\_\_\_\_\_\_ bpm
    \item \textbf{Rhythm}：Atrial - Regular / Irregular \quad Ventricular - Regular / Irregular
    \item \textbf{P Waves}：Present? \_\_\_\_ \quad Appearance? \_\_\_\_\_\_
    \item \textbf{PR Interval}：\_\_\_\_\_\_ 秒 （正常 0.12-0.20 秒）
    \item \textbf{QRS}：\_\_\_\_\_\_ 秒 （正常 $<$0.12 秒）
    \item \textbf{QT Interval}：\_\_\_\_\_\_ 秒 （正常 $<$0.44 秒或 $<$1/2 R-R）
    \item \textbf{ST Segment}：Okay / Elevated / Depressed
    \item \textbf{T Waves}：Normal / Peaked / Inverted / Biphasic
\end{enumerate}

\begin{center}
\fbox{\textbf{Rhythm Interpretation (心律判讀結論)}：\_\_\_\_\_\_\_\_\_\_\_\_\_\_\_\_\_\_\_\_\_\_\_\_\_\_\_\_\_\_\_\_\_\_\_}
\end{center}

\vspace{1cm}

\begin{importantbox}
\textbf{緊急處置提醒}
\begin{itemize}[leftmargin=*]
    \item \textcolor{emergency_red}{\textbf{VF/Pulseless VT}}：立即 CPR + Defibrillation
    \item \textcolor{emergency_red}{\textbf{Asystole/PEA}}：立即 CPR + Epinephrine
    \item \textcolor{emergency_red}{\textbf{Unstable tachycardia}}：Synchronized cardioversion
    \item \textcolor{emergency_red}{\textbf{Symptomatic bradycardia}}：Atropine / Transcutaneous pacing
    \item \textcolor{emergency_red}{\textbf{STEMI}}：立即通知醫師，準備緊急心導管
\end{itemize}
\end{importantbox}

\vfill
\begin{center}
\textcolor{emergency_gray}{\small 本教材僅供教學參考，臨床判斷請依據完整病史與醫師指示}\\
\textcolor{emergency_gray}{\small 原始圖片來源：The Visual Nurse (@thevisualnurse) - www.thevisualnurse.com}
\end{center}

\end{document}
